\documentclass[10pt,a4paper]{amsart}
\usepackage[margin=19mm]{geometry}
\usepackage{amsmath,amssymb,mathtools}
\usepackage[scr=rsfs,cal=euler]{mathalfa}
\usepackage{xfrac}
\usepackage[unicode,hidelinks,bookmarks=false]{hyperref}
\newcommand{\ie}{i.e.\ }
\newcommand{\cf}{cf.\ }
\newcommand{\resp}{respectively\ }
\newcommand{\Subsection}{\subsection}
\providecommand{\nobreakdash}{\nobreak\hbox{-}}
\setlength{\emergencystretch}{2em}
\multlinegap0pt
\usepackage{amssymb}
\usepackage{mathtools}
\usepackage[shortlabels]{enumitem}
\newtheorem{proposition}{Proposition}
\newtheorem{lemma}{Lemma}
\newcommand{\Ric}{\operatorname{Ric}}
\renewcommand{\d}{\partial}
\newcommand{\ddbar}{\sqrt{-1}\d\overline{\d}}
\newcommand{\e}{\mathrm e}
\newcommand{\vep}{\varepsilon}
\title{The top Yau--Yang conjecture for K\"ahler manifolds with positive sectional curvature}
\author{Ved V. Datar, Vamsi P. Pingali,, Harish Seshadri}
\date{Source: arXiv:2606.19806v1 (CC BY 4.0)}
\begin{document}
\maketitle

\noindent\emph{Source and licence.} The top Yau--Yang conjecture for K\"ahler manifolds with positive sectional curvature. Version arXiv:2606.19806v1 (CC BY 4.0), distributed by arXiv under the Creative Commons Attribution 4.0 International licence, http://creativecommons.org/licenses/by/4.0/, as recorded in the arXiv licence metadata for this exact version and shown on the abstract page https://arxiv.org/abs/2606.19806v1. Full licence notice: \texttt{licenses/arxiv-2606.19806-CC-BY-4.0.txt}.\par\smallskip

\noindent\textit{Editorial scope.} Counted unit: the article's lemma lem:mixedterms (source0.tex lines 324--335). The article's complete original proof of it is delivered with it (source0.tex lines 348--448, one \texttt{proof} environment of 101 lines, which the article places away from the statement). No mathematics is written, repaired or re-derived: the statement and the proof are the article's own lines, in the article's own notation, and no retained line is substituted or re-flowed.  Retained uncounted as the source-local context the proof needs: lines 162--176; lines 234--238; lines 316--320.  Nothing was omitted from any retained span, so the package carries no omission record.  No retained line contains a citation, so the wrapper carries no bibliography and no bibliography entry.  No retained line contains a figure environment, an image-inclusion command or drawing code, so no image is delivered and the package declares no asset dependency.  Editorial formatting only, in addition: an AMS \texttt{amsart} preamble that restates the article's own declarations and macros used by the retained lines, namely the environments \texttt{proposition}, \texttt{lemma} on separate counters, together with the standard packages those lines need.  The retained environments print in this document as Proposition 1, Lemma 1 in order of appearance, whereas the article prints the same items with the numbering of its own full document.  The complete native source of the article as distributed in the arXiv e-print for this exact version is bundled unchanged as \texttt{sources/203173/source0.tex}.
\begin{proposition}[$L^2$ canonical section]\label{prop:section}
Let $u_t$, $0\le t\le1$, be the heat-flow smoothing of $\varphi$, with $u_0=\varphi$. There exist $q\gg1$, a non-zero holomorphic section
\[
        s\in H^0(X,K_X),
\]
and a constant $C$ such that
\[
        \int_X \|s\|^2_{q u_t}\,\omega^n\le C,
        \qquad 0\le t\le1.
\]
Here
\[
        \|s\|^2_{q u}:=|s|_\omega^2\e^{-q u}.
\]
\end{proposition}

\begin{equation}\label{eq:BT-convention}
\int_X \eta\, f_{\vep,t}\zeta_{\vep,t}(s)^k\alpha_t^m\omega^l
\longrightarrow
\int_X \eta\, f_\vep\zeta_\vep(s)^k\alpha^m\omega^l
\end{equation}

\begin{equation}\label{ineq:propsofHep}
        0\leq H(h)\leq h,
        \qquad
        h\,\ddbar\log(1+h)\leq\ddbar H(h),
\end{equation}

\begin{lemma}\label{lem:mixedterms}
Fix $\vep>0$.  If $1\leq l\leq n$, $0\leq k,m\leq n-1$, and $k+l+m=n$, then there is a constant $C_{k,l,m}(\vep)$, independent of $a$, such that
\begin{equation}\label{ineq:mixedtermsineqone}
\int_X\chi_a^{4n-2l}f_\vep\zeta_\vep(s)^k\alpha^m\omega^l
\leq C_{k,l,m}(\vep).
\end{equation}
Moreover, for $0\leq p\leq n-1$ there is a constant $D_p(\vep)$, independent of $a$, such that
\begin{equation}\label{ineq:mixedtermsineqtwo}
\int_X\chi_a^{4n}\zeta_\vep(s)^p\alpha^{n-p}
\leq \frac{D_p(\vep)}{a}+q^p\int_X\alpha^n .
\end{equation}
\end{lemma}

\begin{proof}[Proof of Lemma \ref{lem:mixedterms}]
Assume first that \eqref{ineq:mixedtermsineqone} has been proved.  We prove \eqref{ineq:mixedtermsineqtwo} by induction on $p$.  The case $p=0$ is immediate, with $D_0(\vep)=0$.  Suppose the estimate is known at level $p$.  Using regularisation and $V_\vep\leq f_\vep$,
\begin{align*}
\int_X\chi_a^{4n}\zeta_\vep(s)^{p+1}\alpha^{n-p-1}
&=\int_X\chi_a^{4n}\zeta_\vep(s)^p(\ddbar V_\vep+q\alpha)\alpha^{n-p-1}\\
&\leq \frac{C}{a}\int_X\chi_a^{4n-2}f_\vep\zeta_\vep(s)^p\alpha^{n-p-1}\omega
+q\int_X\chi_a^{4n}\zeta_\vep(s)^p\alpha^{n-p}\\
&\leq \frac{D_{p+1}(\vep)}{a}+q^{p+1}\int_X\alpha^n,
\end{align*}
where the first term is absorbed into $D_{p+1}(\vep)/a$ by \eqref{ineq:mixedtermsineqone}.  This proves \eqref{ineq:mixedtermsineqtwo}.

It remains to prove \eqref{ineq:mixedtermsineqone}.  We prove the following regularized estimate, uniform in the smoothing parameter: for $0<t\le1$ and $k+m+l=n$,
\begin{equation}\label{ineq:mixedtermsineqone-reg}
\int_X\chi_a^{4n-2l}f_{\vep,t}\zeta_{\vep,t}(s)^k\alpha_t^m\omega^l
\leq C_{k,l,m}(\vep),
\end{equation}
with $C_{k,l,m}(\vep)$ independent of $a$ and $t$.  Then \eqref{ineq:mixedtermsineqone} follows from \eqref{eq:BT-convention}.

We induct on $j=k+m=n-l$.  When $j=0$,
\[
\int_X\chi_a^{2n}f_{\vep,t}\omega^n
\leq \frac1{\vep^2}\int_X\Vert s\Vert_{q u_t}^2\omega^n
\leq \frac{C}{\vep^2}
\]
by Proposition \ref{prop:section}, uniformly in $a$ and $t$.
Assume \eqref{ineq:mixedtermsineqone-reg} is known for all triples with $k+m\leq j$.  Let $(k,m,l)$ satisfy $k+m=j$ and $k+m+l=n$.  If $l=1$ there is no level $j+1$ estimate to prove, so assume $l\ge2$.

First consider the triple $(k,m+1,l-1)$.  Put
\[
        N:=4n-2l+2,
        \qquad
        T_t:=\zeta_{\vep,t}(s)^k\alpha_t^m\omega^{l-1},
\]
and
\[
        J_{a,t}:=\int_X\chi_a^N f_{\vep,t}\alpha_t\wedge T_t .
\]
Since $\alpha_t=\ddbar u_t$ and $T_t$ is closed, integration by parts and absolute values give
\begin{align}
J_{a,t}
&\leq \frac{C}{a}\int_X\chi_a^{N-2}f_{\vep,t}\zeta_{\vep,t}(s)^k\alpha_t^m\omega^l
+
\left|\int_X\chi_a^N\sqrt{-1}\,\partial f_{\vep,t}\wedge\bar\partial u_t\wedge T_t\right| .
\label{ineq:regularized-J-first}
\end{align}
The first term is bounded by the induction hypothesis.  For the second term, Cauchy--Schwarz gives
\begin{align}
&\left|\int_X\chi_a^N\sqrt{-1}\,\partial f_{\vep,t}\wedge\bar\partial u_t\wedge T_t\right|        \notag\\
&\leq
\left(\int_X\chi_a^N
\frac{\sqrt{-1}\,\partial f_{\vep,t}\wedge\bar\partial f_{\vep,t}}{f_{\vep,t}}
\wedge T_t\right)^{1/2}
\left(\int_X\chi_a^N f_{\vep,t}\sqrt{-1}\,\partial u_t\wedge\bar\partial u_t\wedge T_t\right)^{1/2}.
\label{ineq:regularized-CS}
\end{align}
The quotient is interpreted by its smooth extension across $Z(s)$; locally $f_{\vep,t}=\rho |g|^2$ with $\rho>0$ smooth and $g$ holomorphic.  Since $|\nabla u_t|\leq C$,
$\sqrt{-1}\,\partial u_t\wedge\bar\partial u_t\leq C\omega$, and since $0\leq\chi_a\leq1$,
$\chi_a^N\leq\chi_a^{N-2}=\chi_a^{4n-2l}$.  Hence the second factor in \eqref{ineq:regularized-CS} is bounded by the induction hypothesis, and
\begin{equation}\label{ineq:regularized-CS-reduced}
\left|\int_X\chi_a^N\sqrt{-1}\,\partial f_{\vep,t}\wedge\bar\partial u_t\wedge T_t\right|
\leq
C\left(\int_X\chi_a^N
\frac{\sqrt{-1}\,\partial f_{\vep,t}\wedge\bar\partial f_{\vep,t}}{f_{\vep,t}}
\wedge T_t\right)^{1/2}.
\end{equation}
The Bochner identity for the canonical section $s$ is
\[
\ddbar f_{\vep,t}
=f_{\vep,t}\Ric_\omega-qf_{\vep,t}\alpha_t
+\frac{\sqrt{-1}\,\partial f_{\vep,t}\wedge\bar\partial f_{\vep,t}}{f_{\vep,t}}.
\]
Since $\Ric_\omega\geq0$,
\begin{align}
\int_X\chi_a^N
\frac{\sqrt{-1}\,\partial f_{\vep,t}\wedge\bar\partial f_{\vep,t}}{f_{\vep,t}}\wedge T_t
&\leq
\int_X\chi_a^N\ddbar f_{\vep,t}\wedge T_t+qJ_{a,t}                                      \notag\\
&=\int_X f_{\vep,t}\ddbar(\chi_a^N)\wedge T_t+qJ_{a,t}                                  \notag\\
&\leq C+qJ_{a,t},\label{ineq:regularized-bochner-estimate}
\end{align}
where the last line uses the cutoff estimate and the induction hypothesis.  Combining \eqref{ineq:regularized-J-first}, \eqref{ineq:regularized-CS-reduced}, and \eqref{ineq:regularized-bochner-estimate},
\[
        J_{a,t}\leq C+C(C+qJ_{a,t})^{1/2}\leq C+C\sqrt{J_{a,t}},
\]
so $J_{a,t}\leq C$, uniformly in $a$ and $t$.

Now consider the triple $(k+1,m,l-1)$.  Using \eqref{ineq:propsofHep} at the smooth level,
\begin{align}
&\int_X\chi_a^N f_{\vep,t}\zeta_{\vep,t}(s)^{k+1}\alpha_t^m\omega^{l-1}      \notag\\
&\leq
\int_X\chi_a^N\zeta_{\vep,t}(s)^k\ddbar H(f_{\vep,t})\alpha_t^m\omega^{l-1}
+q\int_X\chi_a^N f_{\vep,t}\zeta_{\vep,t}(s)^k\alpha_t^{m+1}\omega^{l-1}      \notag\\
&\leq\left|\int_X H(f_{\vep,t})\ddbar(\chi_a^N)\wedge\zeta_{\vep,t}(s)^k\alpha_t^m\omega^{l-1}\right|
+q\int_X\chi_a^N f_{\vep,t}\zeta_{\vep,t}(s)^k\alpha_t^{m+1}\omega^{l-1}      \notag\\
&\leq
\frac{C}{a}\int_X\chi_a^{N-2}f_{\vep,t}\zeta_{\vep,t}(s)^k\alpha_t^m\omega^l
+q\int_X\chi_a^N f_{\vep,t}\zeta_{\vep,t}(s)^k\alpha_t^{m+1}\omega^{l-1}
\leq C.\label{ineq:regularized-kplusone}
\end{align}
The last line uses the induction hypothesis for the first term and the already proved $(k,m+1,l-1)$ estimate for the second.  This completes the induction, proves \eqref{ineq:mixedtermsineqone-reg}, and hence proves \eqref{ineq:mixedtermsineqone}.
\end{proof}

\end{document}
